\chapter{相关概念与技术基础}

本章旨在为后续章节的深入论述奠定坚实的理论基础。首先，将详细介绍低地球轨道（LEO）卫星网络与物联网融合的背景，阐明其系统架构、数据特征以及核心挑战。其次，将深入剖析在网计算（In-Network Computing）这一关键使能技术，重点介绍作为其实现基础的P4可编程数据平面技术。最后，将综合前两部分内容，系统性地分析在LEO卫星网络这一特定场景中部署在网计算所面临的独特机遇与严峻挑战，从而为本文的研究工作提供清晰的技术图景。

\section{LEO卫星网络与物联网融合}

\subsection{LEO卫星网络体系架构}

LEO卫星网络作为一种新兴的通信基础设施，其体系架构与传统的地面网络或其他轨道（如GEO、MEO）的卫星系统均有显著不同。一个典型的LEO卫星通信系统，从功能上可以划分为三个相互协作的核心部分：空间段（Space Segment）、地面段（Ground Segment）和用户段（User Segment）。这三部分共同构成了一个完整的、端到端的通信链路，为全球用户提供服务。

空间段是LEO网络的核心，由部署在近地轨道（通常高度在500公里至2000公里之间）的成百上千颗卫星组成的星座构成。与看起来静止不动的GEO卫星不同，LEO卫星以极高的速度（约7.6 km/s）环绕地球运行，单颗卫星对地面某点的可见时间仅为几分钟。为了提供连续覆盖，必须依赖大规模的卫星星座。现代LEO星座的一个关键技术是星间链路（ISL），特别是高速、大容量的激光星间链路。通过ISL，相邻的卫星可以直接通信，形成一个在太空中动态连接的网状网络。这种“太空骨干网”允许数据在卫星之间进行多跳转发，极大地降低了对地面中继站的依赖，实现了真正意义上的全球路由和低延迟通信。根据其在轨处理能力，卫星载荷可分为“透明转发”（Transparent Payload）和“再生处理”（Regenerative Payload）两种。透明载荷仅对信号进行频率转换和放大，如同一个简单的“弯管”（Bent-pipe），所有基带处理和路由决策都在地面完成，这是3GPP Release-17 NTN标准初期的主要模式。再生载荷则在卫星上集成了基带处理、解调/编码、交换路由等功能，相当于将一个地面基站或路由器搬到了太空，具备更强的在轨处理能力和灵活性，是未来发展的趋势。

地面段是连接空间段与地面互联网的桥梁，主要包括关口站（Gateway Station）、测控站（Telemetry, Tracking, and Command, TT\&C Station）和网络操作中心（Network Operations Center, NOC）。关口站是大型地面站，通过大口径天线与卫星星座连接，负责将来自卫星的业务流量接入地面核心网（如互联网、移动核心网），以及将地面流量上传至卫星。测控站则负责跟踪卫星的轨道、姿态，发送控制指令，并接收卫星回报的遥测数据，以确保卫星的正常运行。NOC是整个星座的大脑，负责对网络资源进行集中管理、调度和监控，包括路由计算、频谱分配、用户管理等。

用户段涵盖了所有接入LEO卫星网络的终端设备。这些终端的形态多种多样，以适应不同的应用场景。对于宽带互联网接入服务，用户终端通常是小型的、具备自动跟踪卫星能力的相控阵天线，如Starlink的用户终端（Dishy）。对于移动通信，未来的智能手机将能够直接连接卫星（Direct-to-Cell）。而对于本研究聚焦的物联网应用，用户段则由海量的、通常是低功耗、低成本的物联网终端构成，它们可能分布在农业、物流、环境监测等各种场景中，通过卫星实现数据的采集与回传。

为了更具体地理解现代LEO巨型星座的规模与特征，下表对当前主流的几个系统进行了对比。

\begin{table}[htbp]
    \centering
    \caption{主流LEO巨型星座体系架构对比 (数据来源: 公开资料)}
    \begin{tabular}{llll}
        \toprule
        \textbf{特征} & \textbf{SpaceX Starlink} & \textbf{OneWeb} & \textbf{Amazon Kuiper} \\
        \midrule
        \textbf{轨道高度} & \textasciitilde550 km & \textasciitilde1200 km & \textasciitilde590-630 km \\
        \textbf{计划卫星数量} & \textasciitilde12,000+ & \textasciitilde648 (第一阶段) & \textasciitilde3,236 \\
        \textbf{星间链路技术} & 激光链路 (V1.5及以后版本) & 初期无，后续计划 & 激光链路 \\
        \textbf{用户终端带宽 (下行)} & 25-220 Mbps & 高达 195 Mbps & 高达 400 Mbps \\
        \textbf{用户终端带宽 (上行)} & 5-25 Mbps & 高达 32 Mbps & 未公布 \\
        \textbf{在轨处理架构} & 再生处理能力不断增强 & 初期为透明转发 & 具备在轨处理能力 \\
        \bottomrule
    \end{tabular}
    \label{tab:leo_constellations}
\end{table}

此表清晰地揭示了当前LEO网络发展的共同趋势：大规模部署、低轨道运行以保证低延迟、以及普遍采用星间链路构建太空网络。这些架构特征，特别是星间链路的普及和在轨处理能力的引入，为在网络内部署计算任务提供了物理基础和可能性。

\subsection{卫星物联网的数据流特征}

卫星物联网（SIoT）所承载的数据流，在特性上与传统的人与人（Human-to-Human）或人与机器（Human-to-Machine）通信有着本质的区别。准确理解其数据流特征，是设计高效网络协议和处理机制的前提。SIoT的流量模式通常被归类为海量机器类型通信（massive Machine Type Communications, mMTC），其核心特征可以概括为：海量连接、数据突发、包长极短。

首先，\textbf{海量连接}（Massive Connectivity）是其最显著的特征。与服务于人类用户的宽带网络不同，SIoT旨在为数以千万计甚至亿计的设备提供连接。这些设备（如环境传感器、资产追踪器、智能农具）可能在广阔的地理范围内密集部署。如此庞大的终端数量，对网络的接入控制、地址管理和信令系统构成了前所未有的压力。传统的为每个连接进行握手和资源分配的模式，在如此巨大的连接规模下会产生难以承受的信令开销。

其次，\textbf{数据突发}（Sporadic/Bursty Traffic）是其典型的行为模式。大多数物联网设备并非持续不断地发送数据，而是在大部分时间处于休眠或静默状态以节省能量，仅在特定事件触发时（如检测到异常、达到预设时间点）才被激活，并发送一个或一小批数据。这种流量模式导致信道负载在时间和空间上极不均衡：可能长时间空闲，又在某个瞬间出现大量设备同时尝试接入，形成“信令风暴”或“接入拥塞”，对网络的随机接入信道（RACH）造成巨大冲击。

再次，\textbf{包长极短}（Short Packets）是其数据的普遍形态。物联网应用的数据传输通常是为了上报状态、位置或一个简单的测量值，因此数据包的有效载荷非常小，可能只有几十或几百个字节。在这种情况下，传统网络协议（如TCP/IP）的协议头开销占比会非常高，导致传输效率低下。例如，一个40字节的TCP/IP协议头去承载一个20字节的传感器数据，其有效数据传输效率不足34\%。这种“头重脚轻”的现象在SIoT中极为普遍。

这些数据流特征共同决定了SIoT网络面临的核心挑战，并非源于单个用户的高带宽需求，而是源于\textbf{海量、非协调、突发性短包接入所引发的信令与媒体接入控制（MAC）层面的危机}。传统的基于竞争的接入协议，如ALOHA及其变体，虽然机制简单，但在高并发接入时会因严重的包碰撞而导致信道吞吐量急剧下降，即所谓的“拥塞崩溃”。而基于请求/授权的协议，虽然能避免碰撞，但其为每个短包进行资源申请和分配所带来的信令交互和等待时延，对于突发性业务而言又显得过于笨重和低效。因此，SIoT的数据流特性，强烈地指向了一种技术需求：在网络接入的第一个汇聚点，即LEO卫星上，对数据进行有效的预处理。通过聚合多个短包、过滤冗余数据，可以极大地优化后续的传输和处理流程，从根本上缓解由SIoT独特流量模式带来的网络压力。

\subsection{LEO网络的高动态性挑战}

LEO卫星网络最显著也最具挑战性的特征，源于其组成单元——卫星——的永恒运动。与固定在地面的基站或静止在太空的GEO卫星不同，LEO卫星以大约7.6 km/s的惊人速度相对于地面高速移动，由此引发的一系列动态变化，对网络的设计、控制和运行构成了根本性的挑战。

首先是\textbf{网络拓扑的高度动态性}。从地面用户的视角看，一颗提供服务的卫星仅在头顶上空出现几分钟便会飞越地平线，必须切换到下一颗接踵而至的卫星，这个过程称为“切换”（Handover）。这意味着用户与网络接入点之间的连接是短暂且频繁变化的。从网络自身的视角看，卫星之间的相对位置虽然在同一轨道内相对稳定，但不同轨道面之间的卫星相对位置在不断变化，导致星间链路（ISL）需要不断地断开和重建。因此，整个LEO网络的拓扑结构是一个随时间确定性演化但又极其复杂的动态图。

其次是\textbf{链路特性的剧烈变化}。由于卫星与地面终端之间的距离和相对速度在卫星过顶期间持续变化，无线链路的特性也随之波动。传播延迟（Propagation Delay）会先减小后增大，信号强度（Signal Strength）同样如此。尤为突出的是多普勒频移（Doppler Shift），其变化范围远超地面移动通信，需要终端和卫星具备强大的频率补偿能力才能维持通信。这些链路特性的变化要求物理层和链路层协议必须具备高度的自适应能力。

再次是\textbf{周期性的网络重构}。为了优化全网资源、实现负载均衡，LEO星座运营商会进行周期性的全局网络资源调度和路径重构。针对Starlink网络的深入测量研究发现，该网络存在一个全球同步的、以15秒为周期的“重构间隔”（Reconfiguration Interval）。在每个间隔的边界，用户的通信路径（用户终端-卫星-地面站）可能会被重新指派，这会导致可观测到的时延和吞吐量阶跃式变化。这种主动的、周期性的网络重构，进一步加剧了网络状态的动态性，对需要维持稳定连接或状态的应用（如实时游戏、视频会议，以及本文研究的有状态在网计算）提出了严峻的考验。

综上所述，LEO网络的高动态性是其与生俱来的本质属性。它要求网络中的所有协议、算法和应用都必须具备应对快速变化的能力。对于本文所研究的在网计算而言，这意味着任何需要在网络节点上维持状态（如去重表的哈希值、过滤规则的计数器）的应用，都必须设计出能够在拓扑切换和路径重构中保持状态一致性或实现状态快速迁移的机制。这构成了将传统在网计算思想应用于LEO网络时所面临的最核心、最独特的挑战之一。

\section{在网计算 (In-Network Computing)}

\subsection{P4可编程数据平面技术}

在网计算（In-Network Computing, INC）的理念得以实现，关键在于网络数据平面（Data Plane）具备了前所未有的可编程能力。在这一技术演进浪潮中，P4（Programming Protocol-independent Packet Processors）语言及其所代表的架构思想，扮演了核心的推动者角色。P4是一种高级领域特定语言（Domain-Specific Language），其设计初衷是让网络程序员能够精确定义网络设备（如交换机、路由器、网卡）如何处理数据包，从而将数据平面的行为从固化的硬件逻辑中解放出来。

P4的核心思想基于一个抽象的流水线模型，即\textbf{协议无关交换架构}（Protocol-Independent Switch Architecture, PISA）。PISA将一个可编程交换机的数据处理流程逻辑上划分为三个主要的可编程阶段：

\begin{enumerate}
    \item \textbf{可编程解析器（Programmable Parser）}：当一个原始的数据包（一串比特流）进入交换机时，解析器负责识别和提取其中的协议头。与传统交换机只能识别固定的协议（如Ethernet、IP、TCP）不同，P4的解析器是一个可编程的状态机。程序员可以编写P4代码来定义任意协议头的格式（字段、长度），以及它们在数据包中出现的顺序和条件。这使得交换机能够处理自定义协议、新型协议乃至物联网中各种非标准的封装格式，实现了真正的“协议无关性”。
    
    \item \textbf{可编程匹配-动作流水线（Programmable Match-Action Pipeline）}：这是PISA架构的核心，也是执行在网计算的主要场所。解析器提取出的协议头和元数据（Metadata）被送入一个由多个阶段（Stage）组成的流水线。在每个阶段，数据包的字段可以与一个或多个“匹配-动作表”（Match-Action Table）进行匹配。匹配的类型可以是精确匹配（Exact Match）、最长前缀匹配（Longest Prefix Match, LPM）、三态匹配（Ternary Match，支持通配符）等。一旦匹配成功，就会触发一个预先定义的“动作”（Action）。动作是一段由P4代码编写的指令序列，可以对数据包的头部字段进行修改（如修改IP地址、增减TTL）、对元数据进行读写、执行算术逻辑运算、或与交换机内部的状态化组件（如寄存器、计数器）进行交互。多个匹配-动作阶段串联起来，构成了对数据包进行复杂处理的完整逻辑。
    
    \item \textbf{可编程解解析器（Programmable Deparser）}：在经过匹配-动作流水线的处理后，数据包的头部和元数据可能已经发生了改变。解解析器负责根据程序员定义的逻辑，将这些处理过的头部字段重新序列化，组装成一个完整的比特流，准备从出端口发送出去。
\end{enumerate}

P4程序本身是目标无关的，但它需要被一个特定于硬件目标的编译器编译，才能在具体的交换机ASIC、FPGA或软件交换机上运行。编译器负责将P4程序描述的逻辑流水线，映射到硬件实际拥有的物理资源上（如SRAM、TCAM、ALU等）。这种“语言-架构”分离的设计，使得P4程序具备了高度的可移植性。

为了更清晰地理解P4相比于其前身OpenFlow的革命性进步，下表对二者的可编程模型进行了对比。

\begin{table}[htbp]
    \centering
    \caption{OpenFlow与P4可编程模型对比}
    \begin{tabular}{lll}
        \toprule
        \textbf{特性维度} & \textbf{OpenFlow} & \textbf{P4} \\
        \midrule
        \textbf{解析器灵活性} & 固定解析器 & 可编程解析器 \\
        \textbf{协议支持} & 预定义的、有限的协议头集合 & 协议无关，程序员可自定义任意协议 \\
        \textbf{数据平面逻辑} & 固定的“匹配-动作”范式 & 灵活的、可编程的流水线控制流 \\
        \textbf{状态化处理能力} & 有限的计数器 & 可编程的寄存器、计数器、计量器等状态化原语 \\
        \textbf{目标抽象} & 隐式，依赖于厂商对标准版本的支持 & 显式的架构模型（如PISA, PSA），实现软硬件解耦 \\
        \textbf{开发模式} & “自下而上”，依赖标准演进 & “自顶而下”，由应用需求驱动软件定义 \\
        \bottomrule
    \end{tabular}
    \label{tab:openflow_vs_p4}
\end{table}

从表中可以看出，P4的出现并非对OpenFlow的简单改良，而是一场范式革命。它赋予了网络前所未有的表达能力和灵活性。正是这种能够定义任意协议解析逻辑、并能在流水线中进行状态化计算的能力，使得在网络数据平面上实现如数据去重、过滤等复杂的预处理任务成为可能，为本文的研究提供了最关键的技术基石。

\subsection{在网计算的典型应用场景}

得益于P4等数据平面可编程技术的成熟，在网计算（INC）已经从一个前瞻性的学术概念，发展成为在多个领域拥有实际部署和显著成效的实用技术。这些在地面网络中的成功应用，充分证明了其价值，并为将其拓展至卫星网络等新场景提供了宝贵的经验和启示。

\textbf{网络遥测与诊断}是在网计算最成熟、最广泛的应用之一。传统的网络监控方法（如SNMP、NetFlow）存在采样精度低、数据延迟高、对交换机CPU负担重等问题。基于P4的带内网络遥测（INT）技术彻底改变了这一现状。INT允许交换机在数据包转发过程中，将自身的关键状态信息（如入口/出口端口、时间戳、队列占用长度、链路利用率等）作为一个新的“遥测头”直接附加到原始数据包上。当数据包到达目的地或指定的采集点时，这些沿途收集的遥测信息可以被剥离和分析，从而以线速、逐包、全路径的粒度，精确地重构出数据包所经历的网络状态，为网络故障诊断、性能优化和安全分析提供了前所未有的可见性。

\textbf{流量工程与性能优化}是另一大应用领域。在网计算使得网络能够以微秒级的延迟对流量变化做出反应。例如，P4被用来实现高性能的\textbf{负载均衡器}，通过在数据平面直接计算哈希值并将流量分发到不同服务器，其性能远超基于软件的方案。在\textbf{拥塞控制}方面，P4交换机可以精确地监控队列状态，并实现如ECN标记、QCN等高级拥塞信令机制，甚至可以辅助端到端的拥塞控制协议（如HPCC），为数据中心网络提供近乎零丢包的超低延迟传输。此外，各种复杂的\textbf{主动队列管理（AQM）}算法（如CoDel、PIE）也可以在P4中实现，以更精细化的方式管理交换机缓存，提升网络公平性和整体吞吐量。

\textbf{网络安全}领域同样受益于在网计算。传统依赖于旁路设备或CPU进行流量分析的安全方案，在面对Tbps级别的DDoS攻击时往往力不从心。而基于P4的方案可以将DDoS攻击的特征匹配和过滤逻辑直接部署在网络入口的交换机数据平面上，以线速识别和丢弃攻击流量，在攻击流量对网络内部造成危害之前将其扼杀。同样，可编程防火墙、入侵检测系统（IDS）的规则也可以在P4中实现，大大提升了安全策略的执行效率和灵活性。

\textbf{应用加速与功能卸载}是在网计算的前沿探索方向。研究人员已经成功地将一些原本在服务器上运行的应用功能卸载到P4交换机中。典型的例子包括分布式\textbf{键值存储（Key-Value Store）}的缓存和请求路由、分布式系统中的\textbf{共识协议（如Paxos）}的协调者节点、以及\textbf{机器学习模型推理}等。通过将这些功能的一部分或全部在网络中完成，可以极大地减少服务器的负载，并降低应用访问数据的端到端延迟。

这些成功的应用场景共同揭示了在网计算的核心价值：通过将计算任务智能地迁移到数据平面，可以获得传统计算架构无法比拟的超高性能和超低延迟，从而解决一系列关键的网络和系统问题。

\section{LEO卫星网络中的在网计算：机遇与挑战}

将业已在地面网络中证明其价值的在网计算技术，推广应用于新兴的LEO卫星网络，既充满了诱人的机遇，也面临着前所未有的挑战。这是一个将成熟技术应用于全新环境的交叉研究领域，其独特性值得深入剖析。

\subsection{机遇：结构化拓扑、统一管控、链路带宽瓶颈}

尽管LEO卫星网络环境恶劣、资源受限，但其某些固有特性反而为部署在网计算创造了比地面互联网更为理想的“温床”。

首先，\textbf{链路带宽瓶颈是部署在网计算的最强驱动力}。如前文所述，星地链路是整个系统中最宝贵、最狭窄的瓶颈资源。相比于地面数据中心动辄100 Gbps甚至更高的内部带宽，卫星网络的带宽极其有限。这种极端的资源稀缺性，为任何能够有效减少回传数据量的技术（如在网数据预处理）提供了最充分、最直接的应用动机。在地面网络中，在网计算带来的性能提升可能只是“锦上添花”，但在卫星网络中，它可能是决定系统可用性和经济性的“雪中送炭”的关键。

其次，\textbf{高度结构化的网络拓扑和可预测的动态性}为协同计算提供了便利。与地面互联网那种由无数自治域（AS）组成的、拓扑复杂且难以预测的结构不同，LEO卫星星座的拓扑是高度结构化的。在理想情况下，同一轨道平面内的卫星形成稳定的链状结构，而不同轨道面之间则构成规则的网格状连接，整体呈现出类似环面（Torus）的拓扑特征。更重要的是，这种拓扑的动态演化是基于开普勒轨道力学，完全是确定性和可预测的。这意味着网络控制中心可以提前数小时甚至数天精确计算出未来任意时刻的网络拓扑、星间链路的可用性以及卫星与地面的可见性窗口。这种可预测性为设计主动的、前瞻性的资源调度和状态管理策略提供了可能，这是在随机性更强的地面网络中难以实现的。

再次，\textbf{统一的建设和管控主体}极大地降低了部署和管理的复杂度。一个LEO巨型星座，如Starlink或Kuiper，通常由单一的实体建设、拥有和运营。这意味着运营商对网络中的每一颗卫星（即每一个网络节点）都拥有完全的控制权。这种“绿色地带”（Greenfield）和中央集权的特性，使得在全网范围内部署统一的P4程序、下发一致的控制平面策略、以及进行集中的软件升级变得轻而易举。相比之下，在地面互联网上推广一种新的网络协议或功能，需要获得无数个不同运营商和设备商的共识与协作，过程漫长而艰难。

最后，\textbf{海量闲置资源的潜在价值}。由于地球表面大部分是人烟稀少的海洋、沙漠和极地，LEO星座在任意时刻都有大量的卫星正运行在这些区域上空，其服务负载极低，拥有大量闲置的计算和功率资源。这部分闲置资源构成了一个巨大的、尚未被开发的在轨计算资源池。通过在网计算，可以将一些非实时的、可延迟的计算任务（如科学数据分析、模型训练等）调度到这些“空闲”的卫星上执行，从而将原本被浪费的资源转化为有价值的计算服务，实现“变废为宝”。

\subsection{挑战：动态状态管理、有限计算与存储、功耗与可靠性}

将实验室中的在网计算原型搬到数千公里外的太空，需要克服一系列源于空间环境和卫星平台固有属性的严峻挑战。这些挑战构成了本研究需要攻克的核心技术难关。

\textbf{动态状态管理}是首当其冲且最为独特的挑战。许多有价值的在网计算应用（如本文关注的数据去重、流量聚合统计等）都是有状态的，即它们需要在交换机内部维护一些信息（如哈希表、计数器）来指导后续数据包的处理。在拓扑稳定的地面网络中，状态管理相对简单。但在LEO网络中，由于卫星的快速移动和周期性的网络重构，数据流的路径在不断变化。一个用户终端在几分钟内就可能从卫星A切换到卫星B。这就带来了一个根本性问题：当数据流路径改变时，其在卫星A上积累的状态该如何处理？如果直接丢弃，会导致处理逻辑中断，例如去重失效；如果需要迁移到卫星B，如何设计一个轻量级、低延迟的状态迁移协议，使其开销不至于抵消在网计算带来的收益？这种在毫秒级转发路径上管理分钟级动态网络中的状态，是LEO在网计算面临的核心难题。

\textbf{有限的计算与存储资源}构成了硬性约束。如2.1.3节所述，受SWaP限制，星上P4交换机的硬件资源极其宝贵。可编程流水线的级数、匹配-动作表的总容量（TCAM/SRAM大小）、寄存器阵列的存储空间等都非常有限。一个在地面交换机上可以轻松运行的复杂P4程序，可能会因为资源需求超出了星上芯片的容量而无法编译或加载。因此，为卫星设计的在网计算应用必须在算法层面做到极致的简化和优化，用最少的表项、最少的寄存器和最简单的动作逻辑，来实现所需的功能。

\textbf{严格的功耗与热管理限制}。计算即能耗。在网计算任务的执行会增加交换芯片的活动和功耗，进而产生更多的热量。在卫星这个能量预算紧张、且只能通过辐射进行被动散热的封闭系统中，任何额外的功耗和热量都必须被精确计算和严格管理。一个设计不当的在网计算应用，可能会因为功耗超标而影响到卫星其他关键子系统（如通信载荷）的正常工作，甚至威胁到卫星的生存。

\textbf{极端环境下的可靠性与容错性要求}。太空是一个充满高能粒子辐射的恶劣环境，可能导致芯片发生“单粒子翻转”（SEU）等软错误，造成内存中数据位的改变或计算逻辑的错误。地面设备出现故障可以重启或更换，但在轨卫星的物理维护几乎是不可能的。因此，部署在卫星上的任何软件，包括P4程序和控制平面应用，都必须具备极高的可靠性和鲁棒性。程序中一个微小的bug或逻辑漏洞，都可能导致价值数百万美元的卫星资产功能异常甚至失效。因此，在网计算方案必须包含完善的故障检测、隔离和恢复机制，并经过极其严格的验证和测试，才能考虑在轨部署。

\section{本章小结}

本章系统地梳理了支撑本论文研究的技术背景与核心概念。首先，通过对LEO卫星网络体系架构、卫星物联网数据流特征以及网络高动态性挑战的分析，清晰地描绘出本研究的应用场景及其内在矛盾：一个旨在连接海量、突发、短包物联网设备的全球网络，正面临着由自身架构的高动态性和资源受限性所带来的巨大压力。其次，通过对在网计算与P4可编程数据平面技术的介绍，明确了解决上述矛盾的关键技术手段：利用数据平面的可编程能力，在数据转发路径上以线速进行智能处理。最后，本章深入辨析了在LEO卫星网络中应用在网计算的机遇与挑战。一方面，链路瓶颈、结构化拓扑和统一管控为在网计算提供了绝佳的应用舞台；另一方面，动态状态管理、资源限制、功耗与可靠性等问题，也构成了必须逾越的技术鸿沟。本章的论述为后续章节提出具体的解决方案——即一个面向LEO卫星物联网的在网数据预处理框架——构建了必要的理论基础和问题域界定。

%========================================================
% 第三章：研究点 1
% 面向大规模低轨卫星网络的 LFA（链路洪泛攻击）在网检测与缓解
%========================================================
